% archon:covers AlgebraicJacobian/RiemannRoch/H1Vanishing.lean

\chapter{Vanishing of \(H^1\) for skyscraper sheaves on a curve (RR.2.H\(^1\))}
\label{chap:RR_H1Vanishing}

% SOURCE: Hartshorne, \emph{Algebraic Geometry}, II.1, Exercise~1.16
% (flasque sheaves), Exercise~1.17 (skyscraper sheaves); III.2,
% Proposition~2.5 (flasque sheaves are acyclic) (read from
% references/hartshorne-algebraic-geometry.pdf, PDF pages 67--68 and
% 225). All verbatim quotes pasted in the declaration blocks below were
% copied character-by-character from the rendered page images of the
% scanned PDF (no text layer is present).

\section{Setup and motivation}

This chapter is the project-side build of the
\texttt{RR.2.H\(^1\)} sub-phase, committed in iter-188 of the strategy
revision: \emph{the first cohomology of a closed-point skyscraper sheaf
on a smooth proper curve vanishes}. Concretely, the chapter pins the
identity
\[
  \dim_{\bar k}\, H^1\bigl(C,\, \mathtt{skyscraperSheaf}\;P\;
    (\Module.\mathtt{of}\;\bar k\;\bar k)\bigr) \;=\; 0
\]
on a smooth proper geometrically irreducible curve \(C / \bar k\), where
\(P \in C\) is a closed point and the cohomology is computed through the
project's \(\Module \bar k\)-flavoured \(\Ext\)-cohomology pipeline
\(\HModule \bar k\) (\cref{chap:Cohomology_StructureSheafModuleK}, sibling
\texttt{Cohomology\_StructureSheafModuleK.tex}).

\paragraph{Why a dedicated chapter.} Mathlib (snapshot
\texttt{b80f227}) does \emph{not} ship a flasque-cohomology vanishing
lemma at the \(\Ext\)-cohomology level used by
\(\HModule \bar k\), nor a Serre-duality theorem for curves. The
genus-\(0\) critical path nevertheless needs the skyscraper
\(H^1\)-vanishing twice: once inside the Hartshorne~IV.1.3
\(\chi\)-additivity assembly
(\cref{chap:RiemannRoch_RRFormula}, the parent lemma
\texttt{eulerCharacteristic\_skyscraperSheaf}), and once inside the
\texttt{RR.3} chain that produces a non-constant section
(\cref{chap:RiemannRoch_OCofP}, the downstream consumer
\texttt{h1\_vanishing\_genusZero}). The present chapter unifies the
substrate behind both consumers into a single short, self-contained
project-side build.

\paragraph{Strategy of the chapter.} The argument is the classical
Hartshorne~III~\S2 one: an injective resolution of a flasque sheaf
yields a flasque quotient at each stage (since both the injective object
and the original sheaf are flasque), and the global-sections functor is
exact on a short exact sequence of flasque sheaves. Iterating, every
derived functor of \(\Gamma\) vanishes on a flasque input in positive
degree. The chapter then specialises to the closed-point skyscraper
sheaf, which is flasque by inspection (it is the pushforward, along the
inclusion of a one-point closed subspace, of the constant sheaf on an
irreducible base — and pushforwards of flasque sheaves are flasque).

\paragraph{Standing hypotheses.} Throughout this chapter, \(\bar k\) is
an algebraically closed field and \(C\) is a smooth proper geometrically
irreducible curve over \(\bar k\). The Lean signatures (planned location
\texttt{AlgebraicJacobian/RiemannRoch/H1Vanishing.lean}) carry the same
typeclass package as the rest of the project's curve API
(\(\mathtt{IsProper}\), \(\mathtt{SmoothOfRelativeDimension}\, 1\),
\(\mathtt{GeometricallyIrreducible}\), \(\mathtt{IsIntegral}\)). The
flasque-vanishing lemma of \S2 below is stated more generally for a
noetherian topological space (the natural generality of
Hartshorne~III.2.5); the curve specialisation is read off by instance
synthesis on the underlying space \(C.\mathtt{left.toTopCat}\).

\section{Flasque sheaves and their cohomology vanishing}
\label{sec:H1Vanishing_flasque}

\begin{definition}

\leanok
[Flasque sheaf]
  \label{def:isFlasque_sheaf}
  \lean{AlgebraicGeometry.Scheme.IsFlasque}
  % SOURCE: [Hartshorne], II.1, Exercise~1.16 (definition of flasque
  % sheaf, opening paragraph) (read from
  % references/hartshorne-algebraic-geometry.pdf, PDF page 67); III.2,
  % paragraph preceding Lemma~2.4 (Hartshorne's own recall of the
  % definition) (PDF page 224).
  % SOURCE QUOTE: "Flasque Sheaves. A sheaf \(\mathscr{F}\) on a
  % topological space \(X\) is \emph{flasque} if for every inclusion
  % \(V \subseteq U\) of open sets, the restriction map
  % \(\mathscr{F}(U) \to \mathscr{F}(V)\) is surjective."
  % SOURCE QUOTE: "Recall (II, Ex. 1.16) that a sheaf \(\mathscr{F}\) on
  % a topological space \(X\) is \emph{flasque} if for every inclusion
  % of open sets \(V \subseteq U\), the restriction map
  % \(\mathscr{F}(U) \to \mathscr{F}(V)\) is surjective."
  \textit{Source: Hartshorne, II.1, Exercise~1.16 (definition); recalled
  III.2 preceding Lemma~2.4.}
  Let \((X, \mathcal O_X)\) be a ringed space and let \(\mathcal F\) be a
  sheaf of \(\mathcal O_X\)-modules on \(X\). We say that \(\mathcal F\)
  is \emph{flasque} if, for every inclusion of open sets
  \(V \subseteq U\) in \(X\), the restriction map
  \(\mathcal F(U) \to \mathcal F(V)\) is surjective.

  In the project's \(\Module \bar k\)-flavoured setup the predicate is
  applied to a sheaf
  \(\mathcal F : \mathrm{Sheaf}\;J\;(\Module \bar k)\), where \(J\) is the
  Grothendieck topology of opens on the underlying topological space of
  the curve \(C\). The surjectivity condition is read on the underlying
  presheaf of \(\bar k\)-modules; the project predicate
  \texttt{Scheme.IsFlasque} packages exactly this condition.

  \paragraph{Lean signature scope.} The Lean predicate
  \texttt{AlgebraicGeometry.Scheme.IsFlasque} takes a sheaf
  \(\mathcal F : \mathrm{Sheaf}\;(\Opens.\mathtt{grothendieckTopology}\;X)\;
  (\Module \bar k)\) and asserts that for every pair \(V \le U\) of opens
  of \(X\), the restriction morphism
  \((\mathcal F.\mathtt{val.map}\;(\mathrm{leOp}\;V\;U)) :
  \mathcal F.\mathtt{val.obj}\;(\mathrm{op}\;U) \to
  \mathcal F.\mathtt{val.obj}\;(\mathrm{op}\;V)\) is surjective as a map of
  \(\bar k\)-modules.
\end{definition}

\begin{lemma}

\leanok
[Pushforward of a flasque sheaf is flasque]
  \label{lem:isFlasque_pushforward}
  \lean{AlgebraicGeometry.Scheme.IsFlasque.pushforward}
  \uses{def:isFlasque_sheaf}
  % SOURCE: [Hartshorne], II.1, Exercise~1.16(d) (read from
  % references/hartshorne-algebraic-geometry.pdf, PDF page 67).
  % SOURCE QUOTE: "If \(f: X \to Y\) is a continuous map, and if
  % \(\mathscr{F}\) is a flasque sheaf on \(X\), then \(f_*\mathscr{F}\)
  % is a flasque sheaf on \(Y\)."
  \textit{Source: Hartshorne, II.1, Exercise~1.16(d).}
  Let \(f \colon X \to Y\) be a continuous map of topological spaces and
  let \(\mathcal F\) be a flasque sheaf on \(X\). Then the pushforward
  sheaf \(f_*\mathcal F\) is flasque on \(Y\).
\end{lemma}

\begin{proof}
  \uses{def:isFlasque_sheaf}
  \leanok
  For an inclusion of opens \(V \subseteq U\) in \(Y\), the restriction
  map of \(f_*\mathcal F\) is by definition the restriction map of
  \(\mathcal F\) on the open sets \(f^{-1}(V) \subseteq f^{-1}(U)\) of
  \(X\):
  \[
    (f_*\mathcal F)(U) \;=\; \mathcal F(f^{-1}(U))
    \;\xrightarrow{\;\mathcal F(f^{-1}(V) \subseteq f^{-1}(U))\;}\;
    \mathcal F(f^{-1}(V)) \;=\; (f_*\mathcal F)(V).
  \]
  Since \(\mathcal F\) is flasque on \(X\), the right-hand restriction is
  surjective, hence so is the left-hand one. Therefore \(f_*\mathcal F\)
  is flasque on \(Y\).
\end{proof}

\begin{lemma}

\leanok
[Constant sheaf on an irreducible space is flasque]
  \label{lem:isFlasque_constant_irreducible}
  \lean{AlgebraicGeometry.Scheme.IsFlasque.constant_of_irreducible}
  \uses{def:isFlasque_sheaf}
  % SOURCE: [Hartshorne], II.1, Exercise~1.16(a) (read from
  % references/hartshorne-algebraic-geometry.pdf, PDF page 67).
  % SOURCE QUOTE: "Show that a constant sheaf on an irreducible
  % topological space is flasque. See (I, \S1) for irreducible
  % topological spaces."
  \textit{Source: Hartshorne, II.1, Exercise~1.16(a).}
  Let \(X\) be an irreducible topological space and let \(A\) be a
  \(\bar k\)-module. Then the constant sheaf
  \((\mathrm{constantSheaf}\;J\;\_)\;A\) on \(X\) is flasque.
\end{lemma}

\begin{proof}
  \uses{def:isFlasque_sheaf}
  On an irreducible topological space every non-empty open is connected
  and dense, so the constant presheaf
  \(U \mapsto A\) (with all non-empty restriction maps the identity)
  already satisfies the sheaf condition, and equals its sheafification.
  For every inclusion of non-empty opens \(V \subseteq U\), the
  restriction map is the identity on \(A\) — in particular surjective.
  For an inclusion involving the empty open, the target is the zero
  module, so surjectivity is immediate. Therefore the constant sheaf is
  flasque.

  % NOTE (iter-197 lean-vs-blueprint-checker M-1): The "restriction map
  % is the identity on A" sentence above describes the underlying
  % mathematical fact correctly, but presents the non-empty branch as
  % immediate. In Lean (snapshot `b80f227`) this branch is non-trivial:
  % the constant sheaf
  % `(constantSheaf J D).obj A` is built as the sheafification of the
  % literal constant presheaf, and the sheafification unit
  % \(\eta\colon (\mathrm{Functor.const}\;C^{\mathrm{op}})\,.\,\mathrm{obj}\,A
  % \;\to\;((\mathrm{constantSheaf}\;J\;D)\,.\,\mathrm{obj}\,A)\,.\,\mathrm{val}\)
  % is iso on every non-empty open of an irreducible space, but Mathlib
  % does NOT ship this lemma at the present generality.  The Lean file
  % `AlgebraicJacobian/RiemannRoch/H1Vanishing.lean` (around L150-154 of
  % `Scheme.IsFlasque.constant_of_irreducible`) documents the gap; the
  % empty branch is closed axiom-clean, the non-empty branch carries a
  % typed sorry pending one of two closure routes:
  %
  % Route A (preferred — Mathlib upstreaming candidate, smallest project
  % delta). Provide `Full` and `Faithful` instances for the functor
  % `(constantSheaf J D)` on an irreducible space `X` (equivalently:
  % the `Sheaf.IsConstant` framework for irreducible base spaces).
  % Once these instances are available, the non-empty branch of
  % `constant_of_irreducible` reduces to a single rewrite that pushes
  % the restriction map through the unit iso. Estimated LOC envelope:
  % ~100-150 LOC project-side plus a candidate Mathlib PR.
  %
  % Route B (fallback — no Mathlib upstream needed, larger local
  % construction). Build an alternate sheaf `P'` on `X` with
  % `P'(U) = A` for non-empty `U`, `P'(\bot) = 0`, restriction maps
  % the identity wherever defined; show `P'` is a sheaf on `X`
  % (using irreducibility for the gluing condition); exhibit an
  % isomorphism `(constantSheaf J D).obj A \cong \langle P', _\rangle`
  % by universal property of sheafification. Flasqueness is then
  % immediate from the `P'` description. Estimated LOC envelope:
  % ~150-200 LOC standalone.
  %
  % Route A is the recommended target for iter-197+; Route B is the
  % documented fallback if Route A blocks on a Mathlib gap.
\end{proof}

% NOTE (iter-194 reviewer): The companion Lean substrate helper
% `Scheme.IsFlasque.shortExact_app_surjective` (the
% Hartshorne~II.1~Exercise~1.16(b) consequence packaged here as the
% surjectivity hypothesis) was structurally decomposed iter-194 via a
% `forget₂` bridge to `Sheaf J AddCommGrpCat`: 3 axiom-clean substrate
% helpers landed (`sheafCompose_additive`,
% `sheafCompose_preservesZero`, `Scheme.IsFlasque.toAddCommGrpCat`);
% `Mono SAb.f` + `Epi SAb.g` are now proven axiom-clean inline.  Only
% the typeclass `(sheafCompose (forget₂ ModuleCat AddCommGrpCat))%
% .PreservesHomology` remains as the focused residual gap; iter-195
% closure target.
\begin{lemma}

\leanok
[Cokernel of a flasque-by-flasque short exact sequence is flasque]
  \label{lem:flasque_cokernel_short_exact}
  \lean{AlgebraicGeometry.Scheme.IsFlasque.cokernel_of_shortExact_flasque_flasque}
  \uses{def:isFlasque_sheaf}
  % SOURCE: Project-bespoke substrate helper for the
  % Hartshorne~II.1~Exercise~1.16(c) / III.2.5 chain (iter-193 Lane H
  % prover-side construction). The classical reference is Hartshorne
  % II.1, Exercise~1.16(c); see the verbatim quote in the proof of
  % \cref{thm:H1_vanishing_flasque} below.
  \textit{Source: project-bespoke iter-193 substrate (Hartshorne~II.1,
  Exercise~1.16(c) packaged with the sections-surjectivity input from
  II.1, Exercise~1.16(b) as a hypothesis).}
  Let $X$ be a topological space and let
  \[
    0 \to \mathcal F_1 \to \mathcal F_2 \to \mathcal F_3 \to 0
  \]
  be a short exact sequence in
  $\mathrm{Sheaf}\;(\Opens.\mathtt{grothendieckTopology}\;X)\;
  (\Module \bar k)$. Suppose:
  \begin{enumerate}
    \item $\mathcal F_2$ is flasque (in the sense of
      \cref{def:isFlasque_sheaf}); and
    \item the section-level map
      $\mathcal F_2(U) \to \mathcal F_3(U)$ is surjective for every open
      $U \subseteq X$ (the Hartshorne~II.1, Exercise~1.16(b) input,
      packaged as a hypothesis rather than invoked internally — this
      keeps the lemma's axiom-set clean and lets the consumer site
      supply the surjectivity from its own substrate).
  \end{enumerate}
  Then the cokernel sheaf $\mathcal F_3$ is flasque.

  Concretely, given $V \le U$ in $\Opens X$ and a section
  $t \in \mathcal F_3(V)$: lift $t$ to $\tilde t \in \mathcal F_2(V)$
  via hypothesis~(2) at $V$, extend $\tilde t$ to
  $\tilde T \in \mathcal F_2(U)$ via flasqueness of $\mathcal F_2$, and
  push $\tilde T$ through the map $\mathcal F_2 \to \mathcal F_3$ at
  $U$. Naturality of the morphism on opens identifies the restriction
  of the resulting section in $\mathcal F_3(U)$ with the original $t$.
\end{lemma}

\begin{lemma}

\leanok
[$\Ext^{n+1}$ vanishes from $\Ext^n$ vanishing on the rightmost term, with injective middle term]
  \label{lem:ext_succ_zero_of_injective_lower_zero}
  \lean{AlgebraicGeometry.ext_succ_eq_zero_of_injective_of_lower_zero}
  \uses{def:Scheme_HModule}
  % SOURCE: Project-bespoke substrate helper for the
  % Hartshorne~III.2.5 induction (iter-193 Lane H prover-side
  % construction). The underlying mathematical statement is the
  % classical degree-shifting argument that, for an injective middle
  % term in a short exact sequence, the connecting morphism in the
  % long exact $\Ext$ sequence is surjective above the injective
  % dimension. No external source is cited verbatim.
  \textit{Source: project-bespoke iter-193 substrate.}
  Let $\mathcal C$ be an abelian category equipped with a
  \texttt{HasExt} instance, let $X$ be an object of $\mathcal C$, and
  let
  \[
    0 \to S_1 \to S_2 \to S_3 \to 0
  \]
  be a short exact sequence in $\mathcal C$ whose middle term $S_2$ is
  injective. Let $n_0 \ge 1$ be a natural number and assume that
  $\Ext^{n_0}(X, S_3) = 0$ (i.e.\ every element of $\Ext^{n_0}(X, S_3)$
  vanishes). Then $\Ext^{n_0 + 1}(X, S_1) = 0$.

  The proof is the classical degree-shift. Since $S_2$ is injective and
  $n_0 + 1 \ge 2$, the injective-dimension subsingleton instance forces
  $\Ext^{n_0 + 1}(X, S_2)$ to be subsingleton; combined with exactness
  of the covariant $\Ext$ long exact sequence at $\Ext^{n_0 + 1}(X,
  S_1)$, every $x_1 \in \Ext^{n_0 + 1}(X, S_1)$ is in the image of the
  connecting morphism $\delta \colon \Ext^{n_0}(X, S_3) \to
  \Ext^{n_0 + 1}(X, S_1)$. The hypothesis
  $\Ext^{n_0}(X, S_3) = 0$ then gives $x_1 = 0$. The lemma is consumed
  in the $i \ge 2$ inductive case of \cref{thm:H1_vanishing_flasque}.
\end{lemma}

\begin{lemma}\leanok
[Injective objects of $\mathrm{Sheaf}(X, \Module \bar k)$ are flasque]
  \label{lem:isFlasque_injective}
  \lean{AlgebraicGeometry.Scheme.IsFlasque.injective_flasque}
  \uses{def:isFlasque_sheaf, lem:isFlasque_pushforward}
  % SOURCE: [Hartshorne], III.2, Lemma~2.4 (read from
  % references/hartshorne-algebraic-geometry.pdf — the project's local
  % .md transcript does not extend to this passage, so no verbatim quote
  % is pasted in this block per the iter-197 directive; the citation
  % pointer is recorded for downstream retrieval).
  \textit{Source: Hartshorne, III.2, Lemma~2.4.}

  Let \(I\) be an injective object of
  \(\mathrm{Sheaf}\;(\Opens.\mathtt{grothendieckTopology}\;X)\;
  (\Module \bar k)\). Then \(I\) is flasque in the sense of
  \cref{def:isFlasque_sheaf}: for every inclusion \(V \subseteq U\) of
  opens of \(X\), the restriction map
  \(I(U) \to I(V)\) is surjective on underlying \(\bar k\)-modules.

  The classical argument (Hartshorne~III~Lemma~2.4) uses the
  extension-by-zero \(j_!\) functor: for an open inclusion
  \(j_V \colon V \hookrightarrow X\), one has the short exact sequence
  \(0 \to j_{V!}\mathcal O_V \to j_{U!}\mathcal O_U \to
  j_{U!}\mathcal O_U / j_{V!}\mathcal O_V \to 0\), and injectivity of
  \(I\) makes the induced
  \(\mathrm{Hom}(j_{U!}\mathcal O_U, I) \to
    \mathrm{Hom}(j_{V!}\mathcal O_V, I)\) surjective. Translating
  through the \(j_{(-)!} \dashv j_{(-)}^*\) adjunction (or the
  equivalent presheaf form for the constant ring sheaf \(\bar k\))
  identifies the two Hom-sets with \(I(U)\) and \(I(V)\) respectively,
  delivering the surjective restriction map.
\end{lemma}

\begin{proof}
  \uses{def:isFlasque_sheaf, lem:isFlasque_pushforward}
  % NOTE (iter-197 lean-vs-blueprint-checker J-2): Out-of-loop iter-196
  % directive — closure of this lemma is excluded from the Lane H
  % iter-196 scope per the progress-critic scope reduction. The
  % blocking substrate is Mathlib's extension-by-zero functor `j_!`
  % (the open-immersion pushforward-with-compact-support adjoint),
  % which Mathlib snapshot `b80f227` does not ship in the
  % sheaf-of-modules generality the lemma needs. Estimated closure
  % cost: ~100-150 LOC of project-side `j_!` scaffolding OR a Mathlib
  % upstream PR. The Lean file
  % `AlgebraicJacobian/RiemannRoch/H1Vanishing.lean` carries this
  % declaration as a typed sorry at L613-618. This block is pinned
  % here so that `sync_leanok` and the lean-vs-blueprint-checker can
  % track its closure status independently of the headline theorem
  % below (which depends on it transitively via the private
  % auxiliary `HModule_flasque_subsingleton_aux`).
  See the lemma statement above for the proof sketch; the formalisation
  is deferred to a future iteration after a Mathlib \(j_!\)
  construction is available.
\end{proof}

\begin{theorem}

\leanok
[Flasque sheaves have vanishing higher cohomology]
  \label{thm:H1_vanishing_flasque}
  \lean{AlgebraicGeometry.Scheme.HModule_flasque_eq_zero}
  \uses{def:isFlasque_sheaf, lem:isFlasque_injective,
        lem:flasque_cokernel_short_exact,
        lem:ext_succ_zero_of_injective_lower_zero,
        lem:HModule_flasque_subsingleton_aux,
        lem:HModule_injective_finrank_eq_zero}
  % SOURCE: [Hartshorne], III.2, Proposition~2.5 (read from
  % references/hartshorne-algebraic-geometry.pdf, PDF page 225).
  % SOURCE QUOTE: "Proposition 2.5. \emph{If \(\mathscr{F}\) is a
  % flasque sheaf on a topological space \(X\), then
  % \(H^i(X, \mathscr{F}) = 0\) for all \(i > 0\).}"
  \textit{Source: Hartshorne, III.2, Proposition~2.5, p.~208.}
  Let \(X\) be a topological space and let \(\mathcal F\) be a flasque
  sheaf of \(\bar k\)-modules on \(X\), viewed as an object of
  \(\mathrm{Sheaf}\;(\Opens.\mathtt{grothendieckTopology}\;X)\;
  (\Module \bar k)\). Then, for every \(i \ge 1\),
  \[
    \HModule\;\bar k\;\mathcal F\;i \;=\; 0.
  \]
  In particular, \(\dim_{\bar k} \HModule\;\bar k\;\mathcal F\;1 = 0\).
\end{theorem}

% SOURCE QUOTE PROOF: "Proof. Embed \(\mathscr{F}\) in an injective
% object \(\mathscr{I}\) of \(\mathfrak{Ab}(X)\) and let \(\mathscr{G}\)
% be the quotient:
% \(0 \to \mathscr{F} \to \mathscr{I} \to \mathscr{G} \to 0\).
% Then \(\mathscr{F}\) is flasque by hypothesis, \(\mathscr{I}\) is
% flasque by (2.4), and so \(\mathscr{G}\) is flasque by (II, Ex. 1.16c).
% Now since \(\mathscr{F}\) is flasque, we have an exact sequence (II,
% Ex. 1.16b)
% \(0 \to \Gamma(X, \mathscr{F}) \to \Gamma(X, \mathscr{I}) \to
% \Gamma(X, \mathscr{G}) \to 0\).
% On the other hand, since \(\mathscr{I}\) is injective, we have
% \(H^i(X, \mathscr{I}) = 0\) for \(i > 0\) (1.1Ae). Thus from the long
% exact sequence of cohomology, we get \(H^1(X, \mathscr{F}) = 0\) and
% \(H^i(X, \mathscr{F}) \cong H^{i-1}(X, \mathscr{G})\) for each
% \(i \geq 2\). But \(\mathscr{G}\) is also flasque, so by induction on
% \(i\) we get the result."
\begin{proof}
  \uses{def:isFlasque_sheaf}
  \leanok
  We follow Hartshorne's induction. Embed \(\mathcal F\) into an
  injective object \(\mathcal I\) of the abelian category
  \(\mathrm{Sheaf}\;(\Opens.\mathtt{grothendieckTopology}\;X)\;
  (\Module \bar k)\) — such an injective exists by the
  enough-injectives instance for sheaves of \(\bar k\)-modules on opens
  (\cref{chap:Cohomology_StructureSheafModuleK}). Form the quotient
  short exact sequence
  \[
    0 \to \mathcal F \to \mathcal I \to \mathcal G \to 0
  \]
  in the sheaf category.

  \emph{Step 1: the quotient \(\mathcal G\) is flasque.} The hypothesis
  gives that \(\mathcal F\) is flasque. The injective object
  \(\mathcal I\) is flasque by the Hartshorne~III~Lemma~2.4 input:
  any injective \(\mathcal O_X\)-module is flasque, where the
  \(\mathcal O_X\) here is the constant sheaf of rings \(\bar k\) on
  \(X\), so the sheaf category over the constant sheaf of rings is the
  category of sheaves of \(\bar k\)-modules and the lemma applies. From
  the two flasqueness inputs the quotient \(\mathcal G\) inherits
  flasqueness by Hartshorne~II~Exercise~1.16(c): \emph{if
  \(0 \to \mathcal F' \to \mathcal F \to \mathcal F'' \to 0\) is an
  exact sequence and \(\mathcal F', \mathcal F\) are flasque, then
  \(\mathcal F''\) is flasque}.

  \emph{Step 2: the global-sections functor is exact on the chosen
  short exact sequence.} Since the leftmost term \(\mathcal F\) is
  flasque, the long exact sequence of cohomology starts with a short
  exact sequence on global sections
  (Hartshorne~II~Exercise~1.16(b)):
  \[
    0 \to \Gamma(X, \mathcal F) \to \Gamma(X, \mathcal I) \to
    \Gamma(X, \mathcal G) \to 0.
  \]

  \emph{Step 3: the long exact sequence gives the vanishing.}
  The object \(\mathcal I\) is injective in the abelian sheaf category,
  hence its right-derived global-sections cohomology vanishes in
  positive degree:
  \(\HModule\;\bar k\;\mathcal I\;i = 0\) for all \(i \ge 1\). The long
  exact sequence of the short exact sequence above, applied to the
  derived functor \(\HModule\;\bar k\;{-}\;{\bullet}\) of \(\Gamma\),
  reads in degrees \(0\) and \(1\) as
  \[
    \Gamma(X, \mathcal I) \to \Gamma(X, \mathcal G)
    \to \HModule\;\bar k\;\mathcal F\;1
    \to \HModule\;\bar k\;\mathcal I\;1 = 0.
  \]
  The leftmost map is surjective by Step~2, so the connecting morphism
  to \(\HModule\;\bar k\;\mathcal F\;1\) is the zero map, and the kernel
  of \(\HModule\;\bar k\;\mathcal F\;1 \to 0\) is all of
  \(\HModule\;\bar k\;\mathcal F\;1\); hence
  \(\HModule\;\bar k\;\mathcal F\;1 = 0\).

  \emph{Step 4: higher degrees by induction.} For \(i \ge 2\), the long
  exact sequence yields
  \(\HModule\;\bar k\;\mathcal F\;i \cong \HModule\;\bar k\;\mathcal G\;
  (i - 1)\); the inductive hypothesis on the flasque quotient
  \(\mathcal G\) gives the right-hand side is zero. This closes the
  vanishing in all positive degrees. The chapter consumes only the
  degree-\(1\) instance, but the lemma is stated and proved in full
  generality so it can support later sub-builds (e.g.\ Grothendieck
  vanishing on a one-dimensional scheme).

  % NOTE: the proof structure mirrors Hartshorne~III~Proposition~2.5
  % verbatim. The Lean closure is gated on three substrate inputs:
  % (a) the enough-injectives instance for `Sheaf J (ModuleCat kbar)`
  % already provided by chapter
  % `Cohomology_StructureSheafModuleK.tex`; (b) the project-side
  % flasque-quotient lemma (Hartshorne~II~Ex.~1.16(c)) restated for
  % `ModuleCat kbar`-valued sheaves; (c) the long-exact-sequence
  % carrier for `Abelian.Ext.covariantSequence` at the `HModule kbar`
  % level (the same LES carrier consumed by
  % `eulerCharacteristic_shortExact_add` in
  % `RiemannRoch_RRFormula.tex`). All three are project-side ~50-100
  % LOC apiece.
\end{proof}

\section{Skyscraper sheaves are flasque}
\label{sec:H1Vanishing_skyscraperFlasque}

\begin{lemma}

\leanok
[Skyscraper sheaf as pushforward from a closed point]
  \label{lem:skyscraperSheaf_eq_pushforward}
  \lean{AlgebraicGeometry.Scheme.skyscraperSheaf_eq_pushforward_const}
  \uses{lem:skyscraperSheaf_iso_constantSheaf_punit,
        lem:closedPoint_closure_irreducible}
  % SOURCE: [Hartshorne], II.1, Exercise~1.17 (read from
  % references/hartshorne-algebraic-geometry.pdf, PDF page 68).
  % SOURCE QUOTE: "Skyscraper Sheaves. Let \(X\) be a topological space,
  % let \(P\) be a point, and let \(A\) be an abelian group. Define a
  % sheaf \(i_P(A)\) on \(X\) as follows: \(i_P(A)(U) = A\) if \(P \in
  % U\), \(0\) otherwise. \ldots Show that this sheaf could also be
  % described as \(i_*(A)\), where \(A\) denotes the constant sheaf
  % \(A\) on the closed subspace \(\{P\}^-\), and \(i\colon \{P\}^- \to
  % X\) is the inclusion."
  \textit{Source: Hartshorne, II.1, Exercise~1.17.}
  Let \(X\) be a topological space, let \(P \in X\) be a point with
  closure \(\overline{\{P\}}\), and let \(A\) be a \(\bar k\)-module.
  Write \(i \colon \overline{\{P\}} \hookrightarrow X\) for the closed
  embedding. Then the skyscraper sheaf
  \(\mathtt{skyscraperSheaf}\;P\;A\) on \(X\) is isomorphic to the
  pushforward \(i_*\bigl((\mathrm{constantSheaf}\;{-}\;\_)\;A\bigr)\)
  of the constant sheaf with value \(A\) on \(\overline{\{P\}}\).

  % NOTE (iter-197 lean-vs-blueprint-checker M-2): The Lean signature
  % `Scheme.skyscraperSheaf_eq_pushforward_const` (file
  % `AlgebraicJacobian/RiemannRoch/H1Vanishing.lean`, L818) is
  % materially weaker than the mathematical statement displayed above:
  % the Lean conclusion is
  %     `Nonempty (skyscraperSheaf P A ≅
  %        (pushforward _ _).obj ((constantSheaf _ _).obj A))`,
  % i.e.\ the existence of an isomorphism rather than a concrete one.
  % The `Nonempty` wrapper is forced by the use of `Classical.choice`
  % to unpack the inner iso of
  % \cref{lem:skyscraperSheaf_iso_constantSheaf_punit}, which itself
  % carries a typed sorry pending the Mathlib gap discussed there.
  % When that inner iso is constructed explicitly, the `Nonempty`
  % wrapper can be removed and the present declaration upgraded to
  % a direct concrete-iso statement (no consumer change required at
  % the headline `H1_skyscraperSheaf_finrank_eq_zero` because it
  % does not extract iso morphisms from this lemma).
\end{lemma}

\begin{proof}
  \uses{lem:skyscraperSheaf_iso_constantSheaf_punit}
  \leanok
  Both presheaves send an open \(U \subseteq X\) to \(A\) when
  \(P \in U\) (equivalently, when
  \(U \cap \overline{\{P\}}\) is non-empty, since the closed subset
  \(\overline{\{P\}}\) is irreducible and any open meeting it contains
  \(P\)), and to \(0\) when \(P \notin U\) (equivalently,
  \(U \cap \overline{\{P\}} = \varnothing\)). The restriction maps agree
  by inspection: in the \(P \in V \subseteq U\) case both are the
  identity on \(A\); in the \(P \notin V\) case both factor through
  \(0\). The natural identification on opens is a sheaf morphism, and
  on stalks it is the identity (the stalk at any \(Q \neq P\) outside
  \(\overline{\{P\}}\) is zero on both sides; the stalk at any point of
  \(\overline{\{P\}}\) is \(A\) on both sides). Hence the comparison is
  an isomorphism of sheaves.

  In the Lean formalisation this argument is split into two pieces.
  The \emph{outer} step is the presheaf-level equality
  \[
    \mathtt{skyscraperSheaf}\;P\;A \;=\;
      \bigl(i_*\;\bigr)
      \bigl(\mathtt{skyscraperSheaf}\;\PUnit.\mathrm{unit}\;A\bigr)
  \]
  obtained from Mathlib's \texttt{skyscraperPresheaf\_eq\_pushforward}
  via \texttt{ObjectProperty.FullSubcategory.ext} (axiom-clean, no
  Mathlib gap). The \emph{inner} step is the iso
  \(\mathtt{skyscraperSheaf}\;\PUnit.\mathrm{unit}\;A \cong
    (\mathrm{constantSheaf}\;J_{\PUnit})\,.\,\mathrm{obj}\;A\)
  on the one-point space \(\PUnit\); this is the actual blocking
  technical step and is isolated as
  \cref{lem:skyscraperSheaf_iso_constantSheaf_punit} below. Composing
  the two pieces (with the inner iso unpacked via
  \texttt{Classical.choice} from its current \(\mathtt{Nonempty}\)
  wrapper) gives the lemma.
\end{proof}

\begin{lemma}\leanok
[Inner iso: skyscraperSheaf vs.\ constantSheaf on $\PUnit$]
  \label{lem:skyscraperSheaf_iso_constantSheaf_punit}
  \lean{AlgebraicGeometry.Scheme.skyscraperSheaf_iso_constantSheaf_punit}
  \uses{def:isFlasque_sheaf, lem:isFlasque_constant_irreducible}
  \textit{Project-bespoke sub-lemma} (Lean-substrate technical step
  for \cref{lem:skyscraperSheaf_eq_pushforward}; no separate external
  source — the underlying mathematical content is the one-point case
  of Hartshorne~II~Exercise~1.17 quoted in
  \cref{lem:skyscraperSheaf_eq_pushforward}).

  Let \(\PUnit\) denote the one-point topological space, let
  \(J_{\PUnit}\) denote the Grothendieck topology of opens on
  \(\PUnit\), and let \(A\) be a \(\bar k\)-module. Then there is an
  isomorphism of sheaves of \(\bar k\)-modules on \(\PUnit\)
  \[
    \mathtt{skyscraperSheaf}\;(\PUnit.\mathrm{unit})\;A \;\cong\;
      (\mathrm{constantSheaf}\;J_{\PUnit}\;\Module \bar k)\,.\,\mathrm{obj}\;A.
  \]
\end{lemma}

\begin{proof}
  \uses{def:isFlasque_sheaf, lem:isFlasque_constant_irreducible,
        def:alphaConstToSkyPUnit, def:betaSkyToConstPUnit,
        lem:alpha_beta_eq_toSheafify}
  \leanok
  The space \(\PUnit\) has exactly two opens, \(\bot\) (empty) and
  \(\top\) (the whole space, which contains \(\PUnit.\mathrm{unit}\)).
  We build the iso pointwise on these two opens and check naturality
  on the unique non-identity inclusion \(\bot \le \top\).

  \emph{Forward map.} Apply Mathlib's
  \(\mathtt{constantSheafAdj}\;J_{\PUnit}\;\Module \bar k\;
   h_{T\top}\,.\,\mathtt{homEquiv}\,.\,\mathtt{symm}\) to the natural
  identification
  \(\mathtt{skyscraperSheaf}\;(\PUnit.\mathrm{unit})\;A\,.\,\mathrm{val}\,.
    \,\mathrm{obj}\;(\mathrm{op}\;\top) = A\),
  which holds by \(\mathtt{simp}\;[\mathtt{skyscraperPresheaf}]\) since
  \(\PUnit.\mathrm{unit} \in \top\).

  \emph{Inverse map.} Build pointwise on the two opens of \(\PUnit\):
  \begin{itemize}
    \item at \(\mathrm{op}\;\bot\): both sheaves give the terminal
      object in \(\Module \bar k\) (the constant sheaf is the
      sheafification of the constant presheaf, whose value at the
      empty open is the zero module after sheafification; the
      skyscraper sheaf at the empty open is by construction
      \(0\)). Use \(\mathtt{IsTerminal.uniqueUpToIso}\) to identify
      them;
    \item at \(\mathrm{op}\;\top\): use the constantSheaf-unit-iso on
      \(\PUnit\), which is true since the \(\mathtt{constantSheaf}\)
      functor is fully faithful on \(\PUnit\) (this fully-faithful
      instance is the Mathlib gap described below in the route
      choice).
  \end{itemize}

  \emph{Naturality} on the unique non-identity inclusion
  \(\bot \le \top\): both sides land in the terminal object at
  \(\bot\), so commutativity is automatic by \(\mathtt{Subsingleton}\)
  of the terminal object.

  \paragraph{Mathlib gap (route choice).} The inverse map's
  \(\mathrm{op}\;\top\) component needs that the
  \(\mathtt{constantSheaf}\) functor on \(\PUnit\) is fully faithful.
  Two routes close the gap; both are tracked for the iter-197+
  prover:
  \begin{itemize}
    \item \emph{Route A} (preferred — Mathlib upstreaming
      candidate): provide instances
      \texttt{(constantSheaf (Opens.grothendieckTopology PUnit) D)}
      \(.\) \texttt{Full} and
      \texttt{(constantSheaf (Opens.grothendieckTopology PUnit) D)}
      \(.\) \texttt{Faithful}. Likely the smallest project-side
      change once a Mathlib PR for the analogous fact for the
      one-point space is accepted; the inverse map then follows by
      \texttt{Functor.preimage} of the forward map.
    \item \emph{Route B} (fallback — no Mathlib upstream): directly
      construct the iso pointwise on the two opens of \(\PUnit\) as
      above, using \(\mathtt{Sheaf.IsTerminal}\) at \(\bot\) and an
      explicit value extraction at \(\top\) (the constant sheaf at
      the unique non-empty open of an irreducible space is the
      identity on \(A\) after sheafification — same fact as
      \cref{lem:isFlasque_constant_irreducible} for the non-empty
      branch, used here only at the one-point space). Estimated
      LOC: ~50-80 standalone.
  \end{itemize}
\end{proof}

\begin{lemma}

\leanok
[Closure of a closed point is irreducible]
  \label{lem:closedPoint_closure_irreducible}
  \lean{AlgebraicGeometry.Scheme.PrimeDivisor.closure_isIrreducible}

  Let \(C\) be a Noetherian scheme (in particular the smooth proper
  curve \(C / \bar k\) of this chapter), and let \(P \in C\) be a closed
  point. Then the closure \(\overline{\{P\}} = \{P\}\) is irreducible
  (trivially: a singleton is irreducible as a topological space).

  \textit{Project-bespoke ancillary} (immediate set-theoretic fact;
  no external source pinned).
\end{lemma}

\begin{proof}


\leanok
  A singleton topological space is irreducible: its only non-empty open
  is the whole space, which trivially intersects every other non-empty
  open in a non-empty set (in fact in itself). For a closed point
  \(P \in C\) the closure \(\overline{\{P\}}\) is \(\{P\}\), which is
  homeomorphic to a singleton, hence irreducible.
\end{proof}

\begin{lemma}

\leanok
[Closed-point skyscraper sheaf is flasque]
  \label{lem:skyscraperSheaf_isFlasque}
  \lean{AlgebraicGeometry.Scheme.skyscraperSheaf_isFlasque}
  \uses{def:isFlasque_sheaf}

  Let \(C\) be a smooth proper geometrically irreducible curve over
  \(\bar k\) and let \(P \in C\) be a closed point. Then the
  closed-point skyscraper sheaf
  \[
    k(P) \;:=\; \mathtt{skyscraperSheaf}\;P\;
      (\Module.\mathtt{of}\;\bar k\;\bar k)
  \]
  is flasque as a sheaf of \(\bar k\)-modules on \(C\).

  \textit{Project-bespoke packaging} (no separately named external
  source; the constituent inputs are quoted in their own lemma blocks
  above).
\end{lemma}

\begin{proof}
  \uses{def:isFlasque_sheaf}
  \leanok
  % NOTE (iter-197 lean-vs-blueprint-checker J-1): The blueprint
  % proof block originally listed a four-lemma chain
  % (`skyscraperSheaf_eq_pushforward`, `closedPoint_closure_irreducible`,
  % `isFlasque_pushforward`, `isFlasque_constant_irreducible`) in
  % `\uses{...}`. The Lean closure at
  % `AlgebraicJacobian/RiemannRoch/H1Vanishing.lean:903-939` takes a
  % direct route via `skyscraperPresheaf_map` on a `P.point ∈ V` case
  % split (the `P.point ∈ V` branch closes via `eqToHom` being an iso;
  % the `P.point ∉ V` branch closes via the codomain being terminal /
  % `Subsingleton`). None of the four originally-listed lemmas is used
  % in the Lean proof, so the `\uses{...}` set has been pruned to
  % `def:isFlasque_sheaf` only. The mathematical informal sketch below
  % still presents the conceptually clean four-lemma argument because
  % it is shorter and more illuminating; readers wishing to track the
  % actual Lean proof should consult the file directly.
  By
  \cref{lem:skyscraperSheaf_eq_pushforward}, the skyscraper sheaf is
  isomorphic to the pushforward
  \(i_*\bigl((\mathrm{constantSheaf}\;{-}\;\_)\;
  (\Module.\mathtt{of}\;\bar k\;\bar k)\bigr)\) along the inclusion
  \(i \colon \overline{\{P\}} \hookrightarrow C\). The base
  \(\overline{\{P\}} = \{P\}\) is irreducible
  (\cref{lem:closedPoint_closure_irreducible}), so the constant sheaf on
  it is flasque (\cref{lem:isFlasque_constant_irreducible}). Pushforward
  of a flasque sheaf is flasque
  (\cref{lem:isFlasque_pushforward}). Flasqueness transports along the
  isomorphism of \cref{lem:skyscraperSheaf_eq_pushforward}, so
  \(k(P)\) is flasque.
\end{proof}

\section{\(H^1\) of a closed-point skyscraper sheaf vanishes}
\label{sec:H1Vanishing_skyscraperH1}

\begin{theorem}

\leanok
[\(H^1\) of a closed-point skyscraper sheaf has dimension zero]
  \label{lem:H1_skyscraperSheaf_finrank_eq_zero_main}
  \lean{AlgebraicGeometry.Scheme.H1_skyscraperSheaf_finrank_eq_zero}
  \uses{thm:H1_vanishing_flasque, lem:skyscraperSheaf_isFlasque}

  Let \(C\) be a smooth proper geometrically irreducible curve over
  \(\bar k\) and let \(P \in C\) be a closed point. Then
  \[
    \dim_{\bar k}\, H^1\bigl(C,\, \mathtt{skyscraperSheaf}\;P\;
      (\Module.\mathtt{of}\;\bar k\;\bar k)\bigr) \;=\; 0,
  \]
  where the left-hand side is the \(\bar k\)-finrank of
  \(\HModule\;\bar k\;(\mathtt{skyscraperSheaf}\;P\;(\Module.\mathtt{of}\;
  \bar k\;\bar k))\;1\).

  \textit{Project-bespoke assembly} of the two substrate inputs of this
  chapter (\cref{thm:H1_vanishing_flasque} and
  \cref{lem:skyscraperSheaf_isFlasque}). The chapter's external source
  is Hartshorne~III.2 Proposition~2.5 applied to the flasque skyscraper
  sheaf — both inputs are quoted verbatim in their own lemma blocks
  above.
\end{theorem}

\begin{proof}
  \uses{thm:H1_vanishing_flasque, lem:skyscraperSheaf_isFlasque}
  \leanok
  The closed-point skyscraper sheaf
  \(k(P) = \mathtt{skyscraperSheaf}\;P\;(\Module.\mathtt{of}\;\bar k\;
  \bar k)\) is flasque (\cref{lem:skyscraperSheaf_isFlasque}). By the
  flasque-cohomology vanishing theorem
  (\cref{thm:H1_vanishing_flasque}) applied at \(i = 1\) on the
  underlying topological space \(C.\mathtt{left.toTopCat}\),
  \[
    \HModule\;\bar k\;(k(P))\;1 \;=\; 0.
  \]
  Taking \(\Module.\mathtt{finrank}\;\bar k\) of the zero module gives
  the claimed identity \(\dim_{\bar k} H^1(C, k(P)) = 0\).
\end{proof}

\section{Lean substrate helpers}
\label{sec:H1Vanishing_substrate}

This section gives a one-to-one blueprint entry for each auxiliary Lean
declaration that supports the headline results above. The first group
assembles the long-exact-sequence machinery behind the flasque-vanishing
theorem \cref{thm:H1_vanishing_flasque}; the second sets up the
\(\mathrm{forget}_2\) bridge used to derive the
Hartshorne~II~Exercise~1.16(b) sections-surjectivity input; the third
collects the explicit presheaf maps that build the one-point inner iso
\cref{lem:skyscraperSheaf_iso_constantSheaf_punit}.

\subsection{Long-exact-sequence substrate for flasque vanishing}

\begin{definition}[Canonical injective-embedding short complex]
  \label{def:injectiveSES}
  \lean{AlgebraicGeometry.Scheme.injectiveSES}
  \uses{def:Scheme_HModule}
  For a sheaf \(\mathcal F\) of \(\bar k\)-modules on \(X\), the
  canonical injective embedding \(\mathcal F \hookrightarrow
  \mathrm{Injective.under}\,\mathcal F\) and the cokernel projection
  onto its cokernel assemble into the short complex
  \[
    \mathcal F \;\longrightarrow\; \mathrm{Injective.under}\,\mathcal F
    \;\longrightarrow\; \mathrm{coker}\bigl(\mathcal F \hookrightarrow
    \mathrm{Injective.under}\,\mathcal F\bigr).
  \]
  This is the short complex through which the Hartshorne~III.2.5
  induction runs.
\end{definition}

\begin{proof} Proved directly in Lean (a definitional packaging of the
  injective-embedding morphism with the cokernel projection).
\end{proof}

\begin{lemma}[Injective-embedding short complex is short exact]
  \label{lem:injectiveSES_shortExact}
  \lean{AlgebraicGeometry.Scheme.injectiveSES_shortExact}
  \uses{def:injectiveSES}
  The short complex \cref{def:injectiveSES} is short exact: the
  embedding is a monomorphism, the cokernel projection is an
  epimorphism, and the sequence is exact at the middle term.
\end{lemma}

\begin{proof} Proved directly in Lean (the mono/epi/exactness data of
  the canonical injective embedding and its cokernel).
\end{proof}

\begin{lemma}[Degree-one Ext vanishing from Hom-surjectivity with injective middle term]
  \label{lem:ext_one_eq_zero_of_hom_surjective_of_injective}
  \lean{AlgebraicGeometry.ext_one_eq_zero_of_hom_surjective_of_injective}
  \uses{def:Scheme_HModule}
  Let \(\mathcal C\) be an abelian category with enough \(\Ext\), let
  \(X\) be an object, and let \(0 \to S_1 \to S_2 \to S_3 \to 0\) be a
  short exact sequence whose middle term \(S_2\) is injective. If the
  post-composition map \(\mathrm{Hom}(X, S_2) \to \mathrm{Hom}(X, S_3)\)
  is surjective, then \(\Ext^1(X, S_1) = 0\). This is the degree-one
  base case of the long-exact-sequence collapse.
\end{lemma}

\begin{proof} Proved directly in Lean; the structural skeleton of the
  degree-one case of \cref{thm:H1_vanishing_flasque}, via the covariant
  \(\Ext\) long exact sequence and the injectivity of \(S_2\).
\end{proof}

\begin{lemma}[Injective sheaves have vanishing higher cohomology]
  \label{lem:HModule_injective_finrank_eq_zero}
  \lean{AlgebraicGeometry.Scheme.HModule_injective_finrank_eq_zero}
  \uses{def:Scheme_HModule}
  For an injective object \(\mathcal I\) of the \(\bar k\)-module sheaf
  category on \(X\) and any \(i \ge 1\), the derived global-sections
  cohomology \(\HModule\;\bar k\;\mathcal I\;i\) has \(\bar k\)-finrank
  zero. This supplies the \(\Ext^i(-, \mathcal I) = 0\) input of the
  long-exact-sequence chase.
\end{lemma}

\begin{proof} Proved directly in Lean, from the injective-dimension
  subsingleton instance and finrank-of-subsingleton; it is the
  injective-acyclicity input of \cref{thm:H1_vanishing_flasque}.
\end{proof}

\begin{lemma}[Subsingleton carrier for flasque vanishing]
  \label{lem:HModule_flasque_subsingleton_aux}
  \lean{AlgebraicGeometry.Scheme.HModule_flasque_subsingleton_aux}
  \uses{def:isFlasque_sheaf, lem:injectiveSES_shortExact,
        lem:ext_one_eq_zero_of_hom_surjective_of_injective,
        lem:ext_succ_zero_of_injective_lower_zero,
        lem:isFlasque_shortExact_app_surjective,
        lem:flasque_cokernel_short_exact, lem:isFlasque_injective}
  For a flasque sheaf \(\mathcal F\) of \(\bar k\)-modules on \(X\) and
  any \(n\), the cohomology \(\HModule\;\bar k\;\mathcal F\;(n+1)\) is
  subsingleton. This is the strong-induction-on-\(n\) carrier of
  \cref{thm:H1_vanishing_flasque}: the base case applies the degree-one
  collapse, and the inductive step reduces to the flasque quotient sheaf
  at one lower degree.
\end{lemma}

\begin{proof} Proved directly in Lean by induction on \(n\): the base
  case \(n = 0\) invokes
  \cref{lem:ext_one_eq_zero_of_hom_surjective_of_injective} on the
  canonical injective short exact sequence
  (\cref{lem:injectiveSES_shortExact}), with the Hom-surjectivity input
  supplied by \cref{lem:isFlasque_shortExact_app_surjective} at the
  whole space; the inductive step invokes the higher-degree analogue
  (\cref{lem:ext_succ_zero_of_injective_lower_zero}) with the flasque
  quotient produced by \cref{lem:flasque_cokernel_short_exact} and
  \cref{lem:isFlasque_injective}.
\end{proof}

\subsection{The \(\mathrm{forget}_2\) bridge to Mathlib flasqueness}

\begin{definition}[Sheaf-compose of an additive functor is additive]
  \label{def:sheafCompose_additive}
  \lean{AlgebraicGeometry.sheafCompose_additive}
  \uses{def:isFlasque_sheaf}
  For a Grothendieck topology \(J\) and an additive functor \(F\)
  between preadditive categories admitting sheaf-composition along
  \(F\), the induced functor \(\mathrm{sheafCompose}\;J\;F\) on sheaf
  categories is additive.
\end{definition}

\begin{proof} Proved directly in Lean (additivity checked on the
  underlying presheaf level).
\end{proof}

\begin{definition}[Sheaf-compose of an additive functor preserves zero morphisms]
  \label{def:sheafCompose_preservesZero}
  \lean{AlgebraicGeometry.sheafCompose_preservesZero}
  \uses{def:sheafCompose_additive}
  Under the hypotheses of \cref{def:sheafCompose_additive}, the functor
  \(\mathrm{sheafCompose}\;J\;F\) preserves zero morphisms.
\end{definition}

\begin{proof} Proved directly in Lean, as a consequence of additivity
  (\cref{def:sheafCompose_additive}).
\end{proof}

\begin{definition}[Sheaf-compose of a finite-limit-preserving functor preserves finite limits]
  \label{def:sheafCompose_preservesFiniteLimits}
  \lean{AlgebraicGeometry.sheafCompose_preservesFiniteLimits}
  \uses{def:isFlasque_sheaf}
  If \(F\) preserves finite limits, the source and target abelian
  categories have finite limits, and the sheaf categories have finite
  limits, then \(\mathrm{sheafCompose}\;J\;F\) preserves finite limits.
  This is what lets the bridge functor send a short exact sequence to an
  exact short complex.
\end{definition}

\begin{proof} Proved directly in Lean: the canonical factorisation
  through the fully faithful sheaf-to-presheaf functor preserves and
  reflects finite limits.
\end{proof}

\begin{lemma}[Project flasqueness transfers to Mathlib flasqueness via \(\mathrm{forget}_2\)]
  \label{lem:isFlasque_toAddCommGrpCat}
  \lean{AlgebraicGeometry.Scheme.IsFlasque.toAddCommGrpCat}
  \uses{def:isFlasque_sheaf}
  Applying \(\mathrm{sheafCompose}\) along the forgetful functor
  \(\Module\;\bar k \to \mathrm{AddCommGrp}\) to a flasque sheaf
  \(\mathcal F\) (in the sense of \cref{def:isFlasque_sheaf}) yields an
  \(\mathrm{AddCommGrp}\)-valued sheaf that is flasque in Mathlib's
  sense: sectionwise surjectivity of restriction maps becomes
  epimorphy. The underlying section-level functions are unchanged by the
  forgetful functor.
\end{lemma}

\begin{proof} Proved directly in Lean (epi-iff-surjective in
  \(\mathrm{AddCommGrp}\) applied to the sectionwise-surjective
  restriction maps of \cref{def:isFlasque_sheaf}).
\end{proof}

\begin{lemma}[Sections-surjectivity from a flasque leftmost term, Hartshorne~II.1~Ex.~1.16(b)]
  \label{lem:isFlasque_shortExact_app_surjective}
  \lean{AlgebraicGeometry.Scheme.IsFlasque.shortExact_app_surjective}
  \uses{def:isFlasque_sheaf, lem:isFlasque_toAddCommGrpCat,
        def:sheafCompose_additive, def:sheafCompose_preservesZero,
        def:sheafCompose_preservesFiniteLimits}
  For a short exact sequence \(0 \to S_1 \to S_2 \to S_3 \to 0\) of
  sheaves of \(\bar k\)-modules on \(X\) whose leftmost term \(S_1\) is
  flasque, the section-level map \(S_2(U) \to S_3(U)\) is surjective for
  every open \(U \subseteq X\). This is the Hartshorne~II.1~Exercise
  1.16(b) input consumed inside the flasque-vanishing argument.
\end{lemma}

\begin{proof} Proved directly in Lean via the \(\mathrm{forget}_2\)
  bridge: \cref{lem:isFlasque_toAddCommGrpCat} transfers flasqueness to
  the \(\mathrm{AddCommGrp}\) level, the bridge functor sends the short
  exact sequence to an exact short complex
  (\cref{def:sheafCompose_additive},
  \cref{def:sheafCompose_preservesZero},
  \cref{def:sheafCompose_preservesFiniteLimits}), and Mathlib's
  flasque-implies-sections-epi lemma applies; surjectivity transfers
  back since the forgetful functor leaves the underlying maps unchanged.
\end{proof}

\subsection{Explicit presheaf maps for the one-point inner iso}

\begin{definition}[Constant-to-skyscraper presheaf map on \(\PUnit\)]
  \label{def:alphaConstToSkyPUnit}
  \lean{AlgebraicGeometry.alphaConstToSkyPUnit}
  \uses{lem:isFlasque_constant_irreducible}
  On the one-point space \(\PUnit\), the presheaf morphism from the
  constant presheaf with value \(A\) to the skyscraper presheaf at
  \(\PUnit.\mathrm{unit}\) with value \(A\): the identity (via
  \(\mathrm{eqToHom}\)) on the non-empty open \(\top\), and the terminal
  map on the empty open \(\bot\). It is the forward leg of the inner iso
  of \cref{lem:skyscraperSheaf_iso_constantSheaf_punit}.
\end{definition}

\begin{proof} Proved directly in Lean (an explicit presheaf morphism;
  naturality checked on the unique non-identity inclusion).
\end{proof}

\begin{definition}[Skyscraper-to-constant presheaf map on \(\PUnit\)]
  \label{def:betaSkyToConstPUnit}
  \lean{AlgebraicGeometry.betaSkyToConstPUnit}
  \uses{lem:isFlasque_constant_irreducible}
  On the one-point space \(\PUnit\), the presheaf morphism from the
  skyscraper presheaf at \(\PUnit.\mathrm{unit}\) with value \(A\) to the
  underlying presheaf of the constant sheaf with value \(A\): the
  identity composed with the sheafification unit on \(\top\), and the
  terminal map on \(\bot\). It is the inverse leg, modulo
  sheafification, of \cref{lem:skyscraperSheaf_iso_constantSheaf_punit}.
\end{definition}

\begin{proof} Proved directly in Lean (an explicit presheaf morphism;
  naturality checked on the unique non-identity inclusion).
\end{proof}

\begin{lemma}[Composition of the two presheaf maps is the sheafification unit]
  \label{lem:alpha_beta_eq_toSheafify}
  \lean{AlgebraicGeometry.alphaConstToSkyPUnit_comp_betaSkyToConstPUnit_eq_toSheafify}
  \uses{def:alphaConstToSkyPUnit, def:betaSkyToConstPUnit}
  The composite of \cref{def:alphaConstToSkyPUnit} followed by
  \cref{def:betaSkyToConstPUnit} equals the sheafification unit of the
  constant presheaf on \(\PUnit\). This is the identity that makes both
  directions of the inner iso
  \cref{lem:skyscraperSheaf_iso_constantSheaf_punit} commute.
\end{lemma}

\begin{proof} Proved directly in Lean (a pointwise check on the two
  opens of \(\PUnit\), the non-empty branch collapsing the identity
  composite to the sheafification unit).
\end{proof}

\section{Lean encoding and dependencies}
\label{sec:H1Vanishing_lean_encoding}

\paragraph{Target file.} The chapter's declarations are scheduled into a
new Lean file \texttt{AlgebraicJacobian/RiemannRoch/H1Vanishing.lean}
(to be scaffolded by the iter-191+ prover round). The single exception
is the main output \cref{lem:H1_skyscraperSheaf_finrank_eq_zero_main},
whose \(\lean{...}\) pin reuses the already-existing typed-sorry
declaration
\texttt{AlgebraicGeometry.Scheme.H1\_skyscraperSheaf\_finrank\_eq\_zero}
in \texttt{AlgebraicJacobian/RiemannRoch/RRFormula.lean}; the prover
closes that sorry by importing the chapter's flasque-vanishing and
skyscraper-is-flasque helpers from the new file.

\paragraph{Planned declarations.} In the order in which they appear in
this chapter:
\begin{enumerate}
  \item \texttt{AlgebraicGeometry.Scheme.IsFlasque} — predicate on
    \(\mathrm{Sheaf}\;(\Opens.\mathtt{grothendieckTopology}\;X)\;
    (\Module \bar k)\). Asserts surjectivity of every restriction map.
  \item \texttt{AlgebraicGeometry.Scheme.IsFlasque.pushforward} —
    pushforward of a flasque sheaf is flasque
    (\cref{lem:isFlasque_pushforward}).
  \item \texttt{AlgebraicGeometry.Scheme.IsFlasque.constant\_of\_irreducible}
    — constant sheaf on an irreducible space is flasque
    (\cref{lem:isFlasque_constant_irreducible}).
  \item \texttt{AlgebraicGeometry.Scheme.HModule\_flasque\_eq\_zero} —
    flasque sheaves have zero \(\HModule\) in positive degree
    (\cref{thm:H1_vanishing_flasque}). \emph{Substrate input}: the
    chapter's heaviest lemma; ~100--200 LOC over the iter-191+ run.
  \item \texttt{AlgebraicGeometry.Scheme.skyscraperSheaf\_eq\_pushforward\_const}
    — skyscraper sheaf is pushforward of constant
    (\cref{lem:skyscraperSheaf_eq_pushforward}).
  \item \texttt{AlgebraicGeometry.Scheme.PrimeDivisor.closure\_isIrreducible}
    — singleton closure of a closed point is irreducible
    (\cref{lem:closedPoint_closure_irreducible}).
  \item \texttt{AlgebraicGeometry.Scheme.skyscraperSheaf\_isFlasque} —
    closed-point skyscraper sheaf is flasque
    (\cref{lem:skyscraperSheaf_isFlasque}).
\end{enumerate}

\paragraph{Existing pinned target.}
\texttt{AlgebraicGeometry.Scheme.H1\_skyscraperSheaf\_finrank\_eq\_zero}
in \texttt{AlgebraicJacobian/RiemannRoch/RRFormula.lean} (signature
already declared with a Tier-3 typed sorry as of iter-188; the iter-191
prover replaces the sorry by composing
\texttt{HModule\_flasque\_eq\_zero} of \cref{thm:H1_vanishing_flasque}
with \texttt{skyscraperSheaf\_isFlasque} of
\cref{lem:skyscraperSheaf_isFlasque}). The chapter pins this
declaration under \cref{lem:H1_skyscraperSheaf_finrank_eq_zero_main}.

\paragraph{Mathlib hooks.} The chapter depends on the following Mathlib
infrastructure (snapshot \texttt{b80f227}), all of which is already
consumed elsewhere in the project: the enough-injectives instance for
\(\mathrm{Sheaf}\;J\;(\Module \bar k)\) (the same instance that backs
the project's \(\HModule \bar k\) construction in
\cref{chap:Cohomology_StructureSheafModuleK}); the long-exact-sequence
of \(\Ext\) at \(\HModule \bar k\) carrier level
(\texttt{Abelian.Ext.covariantSequence}, the same LES used by the
\(\chi\)-additivity lemma of
\cref{chap:RiemannRoch_RRFormula}); the Mathlib API on
\(\mathtt{TopCat.Presheaf.SkyscraperSheaf}\) and its presheaf
description \(\mathtt{skyscraperPresheaf}\). The chapter does not
introduce any new \texttt{import} beyond what
\texttt{RRFormula.lean} already pulls.

\paragraph{Downstream consumers.} Two consumers depend on the chapter's
main output \cref{lem:H1_skyscraperSheaf_finrank_eq_zero_main}:
\begin{itemize}
  \item \cref{chap:RiemannRoch_RRFormula}'s parent lemma
    \texttt{eulerCharacteristic\_skyscraperSheaf} consumes it to close
    the H\(^1\) half of \(\chi(k(P)) = 1\); this is the Lane H
    substrate input that the iter-188 strategy revision named.
  \item \cref{chap:RiemannRoch_OCofP} (the \texttt{RR.3} chapter) uses
    the same vanishing as part of its long-exact-sequence chain for the
    structural short exact sequence
    \(0 \to \mathcal O_C \to \mathcal O_C([P]) \to k(P) \to 0\);
    combined with the genus-\(0\) input
    \(\dim H^1(C, \mathcal O_C) = 0\) this delivers
    \(\dim H^0(C, \mathcal O_C([P])) = 2\), the non-constant-section
    output of \texttt{RR.3}.
\end{itemize}

\section{Out of scope}

The following items lie outside this chapter and are scheduled into
sibling sub-builds:
\begin{itemize}
  \item \emph{General coherent-sheaf cohomology vanishing.} Higher
    cohomology vanishing for general coherent sheaves on a curve
    requires Grothendieck vanishing on a one-dimensional scheme
    (Hartshorne~III.2.7), which the chapter does \emph{not} use; the
    chapter is narrowly tailored to the flasque/skyscraper case
    \(H^1(C, k(P)) = 0\) that the genus-\(0\) critical path needs.
  \item \emph{Serre duality for curves.} The general-\(g\) Riemann--Roch
    statement \(\ell(D) - \ell(K - D) = \deg(D) + 1 - g\) routes through
    Serre duality \(H^1(C, \mathcal L) \cong H^0(C, \omega_C \otimes
    \mathcal L^{-1})^*\); Mathlib has neither the dualizing sheaf
    \(\omega_C\) on a scheme nor the duality theorem. The chapter
    sidesteps duality entirely by using flasqueness directly (the
    Hartshorne~III~\S2 order of presentation, not the
    Hartshorne~III~\S7 one).
  \item \emph{The \(H^0\) half of \(\chi(k(P)) = 1\).} The
    one-dimensionality identification \(\dim H^0(C, k(P)) = 1\) is a
    sibling substrate lemma of \cref{chap:RiemannRoch_RRFormula},
    established against the project's four-step
    \(\bar k\)-linear-equivalence chain. The present chapter is the
    complementary \(H^1\) half.
\end{itemize}
